\section{Proofs for robustness, loading transfer, and adaptation}\label{app:robust_adaptive}

\subsection{Proof of Theorem~\ref{thm:sg_full}}
\noindent\emph{Roadmap.} Identify the tail inequalities actually used by the upper proofs, then use a Gaussian subfamily for family-wise minimax lower bounds.
\begin{proof}
For every unit net direction, Chernoff applied to \eqref{eq:joint_sg} gives
$\PP(|\ip UE|>t)\le2e^{-t^2/(2v^2)}$.
The finite union bound proving scan coverage and the additive-slack bound proving Lemma~\ref{lem:residual_fit} use no independence between directions. Both therefore hold with $v$ replacing $\tau$. All subsequent diameter and squared-error inequalities are deterministic functions of these events and moments. Orthogonal changes of coordinates preserve \eqref{eq:joint_sg}. The weighted proofs apply unchanged. The unknown-law family includes $E\sim N(0,v^2I)$. Honesty over the family implies honesty over that subexperiment, and the worst-case risk or expected length over the family is at least its Gaussian value. The Gaussian lower bounds give the matching order.
\end{proof}

\subsection{Proof of Theorem~\ref{thm:mom_full}}
\noindent\emph{Roadmap.} Control every median by a binomial tail, obtain a common excess-tail moment, and reuse the ordered-difference argument.
\begin{proof}
For a direction at level $s$, each block average has variance at most $\sigma^2/k_s$. Chebyshev gives a block-error probability at most $\sigma^2/(k_st^2)$. If the median errs by more than $t$, at least half its blocks err. Independence of blocks, followed by a union over their subsets, gives
\[
 \PP(|d_U-\ip UA|>\lambda_s^{\rm mom}+t)
 \le 2^{-b_s}(1+t/\lambda_s^{\rm mom})^{-b_s},\quad t\ge0.
\]
This follows by substituting $\lambda_s^{\rm mom}=4\sigma/\sqrt{k_s}$ into
$2^{b_s}\{\sigma^2/[k_s(\lambda_s^{\rm mom}+t)^2]\}^{b_s/2}$.
There are at most $e^{H_s}$ directions at level $s$. By the definition of $b_s$,
$\sum_s e^{H_s}2^{-b_s}\le\beta$.
Thus all constraints cover the true tensor with probability at least $1-\beta$.

Since $k_s\ge n/(2b_s)$,
$\lambda_s^{\rm mom}\le4\sqrt2\tau\sqrt{b_s}$.
Define
$T_0=\max_U(|d_U-\ip UA|-\lambda_{s(U)}^{\rm mom})_+$.
For $a\ge0$, the function $b\log(1+a/\sqrt b)$ is increasing for $b>0$: its derivative is
$\log(1+u)-u/[2(1+u)]\ge0$, $u=a/\sqrt b$.
Because every $b_s\ge9$, the union bound yields
\begin{equation}\label{eq:mom_excess}
 \PP(T_0>t)\le\beta(1+t/(12\sqrt2\tau))^{-9},\qquad
 \E T_0^2\le\frac{72}{7}\beta\tau^2.
\end{equation}
For the moment, integrate $2t$ times the tail and use
$\int_0^\infty u(1+u)^{-9}du=1/(8\cdot7)$.

For two feasible tensors every directional difference is at most $2\lambda_s^{\rm mom}$. The net comparison, doubled caps and norm bound prove exactly the previous norm-sensitive diameter. Medians are measurable finite order statistics. The extrema of their compact feasible sections are Borel by the same compact-section argument as before, now treating $(d_U)_U$ as the finite observation vector. The smallest-odd-integer rule gives
$b_s\le11+[H_s+\log(p/\beta)]/\log2\le C_{K,\beta}s\ell_p$.
The sum lemma therefore bounds the diameter by $C\mathfrak r_K$.

For estimation, $F$ is continuous and one-Lipschitz in $B$. Selecting the first dense-sequence element with value at most $\inf_jF(B_j)+\tau$ is measurable. Its comparison with $A$ gives directional differences at most $2\lambda_s^{\rm mom}+2T_0+\tau$. Lemma~\ref{lem:residual_fit}'s deterministic argument gives
\[
 |\widehat\psi-\psi(A)|\le U^{\rm mom}
       +Q_2z_K(T_0+\tau/2)/\sqrt p,
 \qquad |\widehat\psi-\psi(A)|\le2h_0.
\]
Equation~\eqref{eq:mom_excess} and the small-$M$ structural bound prove risk $O(\mathfrak r_K^2)$ exactly as in Theorem~\ref{thm:all_norm_mse}. For the lower bound restrict the replicate family to iid Gaussian errors with covariance $\sigma^2I$. Orthogonal transformation in the replicate index separates the sample mean, whose noise level is $\sigma/\sqrt n$, from mean-independent residuals. Thus the full replicate likelihood second moments equal the sample-mean ones. The earlier Gaussian lower bounds apply to any procedure using all replicates, including randomized ones. Finally, $H_s$ is increasing in $s$, so the replicate condition has the stated sufficient order $p$ for fixed $K,\beta$.
\end{proof}

\subsection{Proof of Proposition~\ref{prop:noise_limits}}
\noindent\emph{Roadmap.} Exhibit exact scalar recovery and a correlated direction with inflated variance.
\begin{proof}
If $|A|<v$, then $\operatorname{sign}(A+v\xi)=\xi$. The estimator
$Z-v\operatorname{sign}(Z)$ is identically $A$, proving (i). The Rademacher moment-generating function obeys $\cosh(tv)\le e^{t^2v^2/2}$.
For (ii), take the unit product tensor with every entry $p^{-K/2}$. Its contraction with $E$ is $v\xi p^{K/2}$, with variance $v^2p^K$. Each individual entry still has variance and sub-Gaussian proxy $v^2$.
\end{proof}

\subsection{Proof of Theorem~\ref{thm:approx_sharp}}
\noindent\emph{Roadmap.} Inflate scan constraints for a bounded approximation residual, then combine two nested lower-bound subclasses.
\begin{proof}
Write $A=B_0+R$, $B_0\in\Theta$, $\norm R_F\le\varepsilon_A$; compactness ensures such a decomposition exists. On the Gaussian scan event every residual contraction of $Z-B_0$ is at most $\lambda_s+\varepsilon_A$, so $B_0$ is feasible. Inflate its projected target interval by $\varepsilon_A$ on both sides. If the feasible set is empty use a singleton; this cannot occur on the coverage event. Since $\norm L_F=1$, the enlarged interval covers $\psi(A)$.
For feasible differences the diameter argument gives
\[
 |C|\le U_{K,p}(M;\alpha)
       +\{Q_2z_K/\sqrt p+2\}\varepsilon_A,
\]
using $\min(a,b+c)\le\min(a,b)+c$. It also gives $|C|\le2h_0+2\varepsilon_A$.
Choose between this procedure and the direct Gaussian interval using their deterministic length bounds. This proves the upper order
$\min\{\tau,\mathfrak r_K+\varepsilon_A\}$.

The class contains $\Theta$, giving lower bound $c\mathfrak r_K$ at zero. It also contains $\{tL:|t|\le\varepsilon_A\}$ because $0\in\Theta$. Comparing zero with $tL$ for $t=\min(\varepsilon_A,\tau/2)$ gives total variation at most
$\tfrac12\sqrt{e^{1/4}-1}<1/2$ and target difference $t$. Hence the lower bound is also $c_\alpha\min(\varepsilon_A,\tau)$. Finally
$\min(\tau,r+\varepsilon_A)\le r+\min(\tau,\varepsilon_A)\le2\max\{r,\min(\tau,\varepsilon_A)\}$
for every $r\ge0$. This establishes both criteria up to fixed constants.
\end{proof}

\subsection{Proof of Proposition~\ref{prop:growing_caps}}
\noindent\emph{Roadmap.} Embed sparse diagonal tensors of rank $s$, then compare with the full norm ball.
\begin{proof}
For a uniform size-$s$ support $S$, set
$A_S=b\sum_{i\in S}E_i$ with
\[
 b=\min\{M/\sqrt s,\tau\sqrt{\log(1+p/(4s^2))}\}.
\]
All constrained ranks are $s$, its norm is $b\sqrt s\le M$, and its target is $sb/\sqrt p$. The full tensor inner product is $b^2|S\cap T|$; the hypergeometric factorial-moment inequality gives
$1+\chi^2(Q,P_0)\le\exp\{(s^2/p)(e^{b^2/\tau^2}-1)\}\le e^{1/4}$.
The common-target lemma proves the lower bound for every $s\le r_*$. Maximize over $s$.
If there are no rank restrictions, the line $\{tL:|t|\le M\}$ belongs to the Frobenius ball. The two-point argument with $t=\min(M,\tau/2)$ gives lower order $\min(M,\tau)$. The direct interval and $[-M,M]$, chosen deterministically by radius, give the matching upper bound for every mean in that ball. No asymptotic fixed-cap argument applies to this enlarged class.
\end{proof}

\subsection{Proof of Proposition~\ref{prop:loading_transfer}}
\noindent\emph{Roadmap.} Add an independently calibrated error event, without assuming stochastic independence.
\begin{proof}
On the intersection of $\ip{L_0}A\in C_0$ and $|\ip RE|\le q_R$, adding the second target contribution gives the asserted interval. A union bound controls failure by $\alpha_0+\alpha_1$. The length is the displayed sum. Chernoff or the Gaussian distribution of $\ip RE$ supplies the calibrations. For a prescribed partition, repeat the argument over blocks using the sum of their allocated failures; their centers may be dependent. Selecting among complete constructions from deterministic diameter bounds uses no data and retains its own coverage guarantee. Applying the result with $L_0=0$ also covers every $L$, although it then reduces to the direct interval.
\end{proof}

\subsection{Proof of Theorem~\ref{thm:efficient_regimes}}
\noindent\emph{Roadmap.} Match elementary constructive upper bounds with the existing small-budget lower bound and a product-direction two-point test.
\begin{proof}
In (i), the rate formula has $x\le1$ and is $M/\sqrt p$. The structural radius is $M\sqrt{s_{\rm cap}/p}$, with fixed cap factor. Theorem~\ref{thm:all_norm_rates} gives the matching lower bound.
For (ii), choose the sign of $U$ so that $\ip LU\ge c\norm L_F$ and compare zero with $tU$, $t=\min(M,\tau/2)$. The full likelihood second moment is $e^{t^2/\tau^2}\le e^{1/4}$ and the target difference is at least $ct\norm L_F$. The common-target lemma gives the lower order at zero. Cauchy--Schwarz supplies structural radius $M\norm L_F$, while the direct Gaussian radius is $z_{1-\alpha/2}\tau\norm L_F$. Their deterministic selection controls worst-case length and is polynomial in the number of observed coordinates. The cap class is used only for the lower inclusion; fixed rank caps are unnecessary here.
\end{proof}

\subsection{Proof of Proposition~\ref{prop:convex_loss}}
\noindent\emph{Roadmap.} Find a one-dimensional Gaussian submodel inside every such convex relaxation.
\begin{proof}
The average of $ME_i$, $i=1,\ldots,p$, is $(M/\sqrt p)L$. Both signs and zero belong to the convex hull, hence so does the asserted line segment. Compare zero with $tL$ for $t=\min(M/\sqrt p,\tau/2)$ and apply the same two-point calculation. At $M=\tau\sqrt p$ it forces order $\tau$. The previous minimax rates divided by $\tau$ tend to zero for each fixed $K\ge3$. This argument changes the honesty class and makes no complexity assumption.
\end{proof}

\subsection{Proof of Theorem~\ref{thm:shifted_adaptation}}
\noindent\emph{Roadmap.} Shift the full-data hard mixtures into an unused orthogonal block.
\begin{proof}
Embed the dimension-$m$ rank-one alternatives from Theorem~\ref{thm:all_norm_rates}, with norm budget $M_0$, in coordinates $\{q+1,\ldots,p\}^K$. Denote them by $B$. Their supports are disjoint from $A_0$, so
$\norm{A_0+B}_F^2=\norm{A_0}_F^2+\norm B_F^2\le M^2$.
Unfolding rank subadditivity shows $A_0+B$ is in the outer cap class. Translation of Gaussian likelihoods by $A_0$ leaves their second moments unchanged. All alternatives have a common increment in the normalized $p$-diagonal target equal to $\sqrt{m/p}$ times their dimension-$m$ target. Their total variation from the law at $A_0$ is bounded by a fixed constant less than $1-2\alpha$. Apply the common-target lemma with baseline target $\psi(A_0)$, or translate the interval by that fixed number. This proves the lower bound.
\end{proof}

\subsection{Proof of Theorem~\ref{thm:adaptive_tangent}}
\noindent\emph{Roadmap.} Recover ranks on one spectral event, bound a second-order projection residual, and condition on the independent pilot.
\begin{proof}
For each unfolding use $1/4$-nets of its two unit spheres, of total cardinality at most $9^{p_k+d_k}$. Product telescoping bounds its operator norm by twice the maximum net contraction. Chernoff and a union bound over the $K$ unfoldings give an event $\mathcal E$ of probability at least $1-\delta$ on which every pilot error operator norm is at most $\eta_k$. No independence between unfoldings is used.
The singular-value perturbation inequality
$|\sigma_j(Z^{(1)}_{(k)})-\sigma_j(A_{(k)})|\le\eta_k$
follows from the variational characterization of singular values: for every unit vector the two image norms differ by at most $\eta_k$, and taking the relevant min/max over subspaces preserves this bound. Thus zero singular values produce observed values at most $\eta_k$, while nonzero values exceeding $2\eta_k$ produce observed values greater than $\eta_k$. All selected ranks are exact on $\mathcal E$. If any is zero then $A=0$ and the singleton covers.

For a positive selected rank $r_k$, write $A_{(k)}=U_kS_kV_k^\top$ using its true nonzero singular values. Because $Q_k$ is the leading observed rank-$r_k$ projector,
\[
 \norm{(I-Q_k)A_{(k)}}_{\rm op}
 \le\sigma_{r_k+1}(Z^{(1)}_{(k)})+\eta_k\le2\eta_k.
\]
For full row rank the residual is zero, and the same inequality holds. Right multiplication by $V_kS_k^{-1}$ gives
\[
 \norm{(I-Q_k)U_k}_{\rm op}
 \le\frac{2\eta_k}{\sigma_{r_k}(A_{(k)})}
 \le\frac{2\eta_k}{\sigma_{r_k}(Z^{(1)}_{(k)})-\eta_k}.
\]
Clipping at one yields $d_k^*$. This calculation uses the gap only on the stated separated class, with a strictly positive observed denominator.

Expand the identity into the $2^K$ mutually orthogonal products of $Q_k$ and $I-Q_k$. The projector $T$ retains exactly the terms with at most one complement. Each term discarded by $T$ appears in at least one product $(I-Q_k)(I-Q_\ell)$. Orthogonality therefore implies
\[
 \norm{(I-T)A}_F^2
 \le\sum_{k<\ell}\norm{(I-Q_k)(I-Q_\ell)A}_F^2
 \le M^2\sum_{k<\ell}(d_k^*d_\ell^*)^2.
\]
The last inequality uses $A=(\bigotimes_kU_k)G$ with $\norm G_F=\norm A_F$ and the operator norm of a product of mode maps. Also $\norm{(I-T)A}_F\le M$. Since $T$ is orthogonal, Cauchy--Schwarz bounds
$|\ip L{(I-T)A}|$ by $B_L$.

Condition on the pilot. The vector $TL$ is now deterministic, and independence makes the second error contraction sub-Gaussian with proxy $v_2\norm{TL}_F$. Conditional on any pilot in $\mathcal E$, the projected interval has error at most $\alpha-\delta$. The direct alternative has the same conditional failure bound and the structural alternative never fails. A choice determined by the pilot thus also has conditional failure at most $\alpha-\delta$. Intersecting with a surely valid structural interval preserves coverage; empty intersections may be replaced by a singleton. Integrating and adding $\PP(\mathcal E^c)\le\delta$ proves honesty. SVDs, finite sums, and tensor matrix products give the stated arithmetic complexity.

For Gaussian artificial splitting the two centered errors are jointly Gaussian with zero cross covariance, each of covariance $2\tau^2I$, hence independent. Zero covariance would not suffice for non-Gaussian errors. No minimax length conclusion on the original unseparated class follows from this proof.
\end{proof}
